\documentclass[11pt]{article}

\usepackage[T1]{fontenc}
\usepackage[utf8]{inputenc}
\usepackage[a4paper,margin=27mm]{geometry}
\usepackage{microtype}
\usepackage{amsmath,amssymb,amsthm,mathtools}
\usepackage{newtxtext,newtxmath}
\usepackage{booktabs}
\usepackage[numbers,sort&compress]{natbib}
\usepackage{tikz}
\usetikzlibrary{arrows.meta,positioning}
\usepackage[colorlinks=true,allcolors=blue!55!black]{hyperref}

\newtheorem{theorem}{Theorem}
\newtheorem{lemma}[theorem]{Lemma}
\newcommand{\Tr}{\operatorname{Tr}}
\newcommand{\rank}{\operatorname{rank}}
\newcommand{\ii}{\mathrm{i}}
\newcommand{\cH}{\mathcal{H}}
\newcommand{\cG}{\mathcal{G}}
\newcommand{\cO}{\mathcal{O}}
\newcommand{\ket}[1]{\lvert #1\rangle}
\newcommand{\bra}[1]{\langle #1\rvert}
\newcommand{\proj}[1]{\lvert #1\rangle\!\langle #1\rvert}

\title{Local stability of a two-copy entanglement-of-purification ansatz at the von Neumann boundary}
\author{Xiaojian Zeng\\
\small PBC School of Finance, Tsinghua University, Beijing 100190, China\\
\small \texttt{xiaojianzeng@foxmail.com}}
\date{}

\begin{document}
\maketitle

\begin{abstract}
Entanglement of purification requires an optimization over purifications, and
its additivity at von Neumann order remains open even for classical states.
For a diagonal two-qubit family recently studied at other R\'enyi orders, we
analyze two copies at $r=1/2$ in the minimal $4\times4$ ancillary split.  We
prove, with a computer-assisted interval certificate, that the
product-canonical purification is a strict local minimum modulo local
ancillary unitaries throughout the full complex $U(16)$ flag manifold,
including rank-opening directions.  The proof combines an 81-dimensional
quotient-Hessian bound with a logarithmic tail estimate.  No global
optimality, enlarged-ancilla, or additivity claim is made.
\end{abstract}

\noindent\textbf{Keywords:} entanglement of purification; von Neumann entropy;
purification landscape; interval arithmetic; local minimum

\section{Introduction}

The entanglement of purification (EoP) of a bipartite state $\rho_{AB}$ is
the least entanglement entropy across $AA':BB'$ among all of its
purifications~\cite{Terhal2002}.  Its regularization has operational meanings
in correlation preparation and visible compression~\cite{Terhal2002,Hayashi2006},
but it is not known in general whether the one-copy quantity is additive.
Numerical work on two-qubit Werner states supplied evidence for
nonadditivity~\cite{ChenWinter2012}, without resolving the von Neumann
problem analytically.

Negari and Baghali Khanian recently isolated a simple diagonal two-qubit
family for which the one-copy optimization can be solved exactly at every
R\'enyi order~\cite{NegariKhanian2026}.  For two copies they proved strict
nonadditivity for every order below one using a signed-canonical witness, and
proved strong additivity, under their stated classical premise, for orders at
least two.  The interval $1\leq\alpha<2$, including the von Neumann point
$\alpha=1$, remains open in that work.  Their finite sign search does not
cover continuous complex purifying flags, while a positive ordinary Hessian
alone cannot settle directions that open zero eigenvalues of the reduced
state.

Here we resolve one deliberately local part of that open landscape.  At the
exact state parameter $r=1/2$, and only in the minimal ancillary split
$(d_{A'},d_{B'})=(4,4)$, we show that the product-canonical two-copy
purification is a strict local minimum modulo local ancillary unitaries over
the complete complex flag manifold $U(16)$.  The tangent space separates
into 112 rank-preserving and 144 support-opening directions.  Removing the
31-dimensional ancillary gauge orbit leaves an 81-dimensional analytic
Hessian.  An outward-rounded ball-arithmetic calculation certifies its least
eigenvalue to exceed $0.159$.  A separate Schur-complement argument uses the
$-t^2\log t^2$ entropy generated at the rank boundary to make this Hessian
result uniform for mixed rank-preserving and support-opening paths.

The conclusion is local and dimension-specific.  It neither minimizes over
distant points of $U(16)$ nor compares larger ancillary splits, and therefore
does not prove two-copy additivity.  Its role is instead to identify a
rigorous obstruction: within the minimal split, a von Neumann improvement at
this state cannot bifurcate infinitesimally from the canonical purification.

\section{Model and main result}

For $0\leq r\leq1$, consider the classical state
\begin{equation}
 \rho_r=\frac{r^2\proj{00}+\proj{01}+\proj{10}+\proj{11}}{3+r^2}.
 \label{eq:state}
\end{equation}
All logarithms below have base two, so entropies are measured in bits.  The
von Neumann EoP is
\begin{equation}
 E_P(\rho_{AB})=\inf_{\ket{\Psi}:\,\Tr_{A'B'}\proj{\Psi}=\rho_{AB}}
 S\!\left(\Tr_{BB'}\proj{\Psi}\right),\qquad
 S(\sigma)=-\Tr(\sigma\log_2\sigma).
 \label{eq:eop}
\end{equation}

For two copies, write $I=(i_1,i_2)$ and $J=(j_1,j_2)$, each ordered
lexicographically as $00,01,10,11$, and let
$P_{IJ}=p_{i_1j_1}p_{i_2j_2}$.  At $r>0$ all sixteen probabilities are
positive.  In the minimal split $\cH_{A'}\otimes\cH_{B'}\cong
\mathbb C^4\otimes\mathbb C^4$, every orthonormal flag family is an
orthonormal basis and hence can be represented by $W\in U(16)$:
\begin{equation}
 \ket{\Psi_W}=\sum_{I,J}\sqrt{P_{IJ}}\,\ket{I}_A\ket{J}_B
 W\ket{IJ}_{A'B'}.
 \label{eq:flags}
\end{equation}
Let
\begin{equation}
 F(W)=S\!\left(\Tr_{BB'}\proj{\Psi_W}\right).
 \label{eq:objective}
\end{equation}
The product-canonical purification is $W_0=I_{16}$.

At $r=1/2$, the normalized one-copy coefficient matrix across $AA':BB'$ is
\begin{equation}
 C_1=\frac{1}{\sqrt{13}}
 \begin{pmatrix}1&2\\2&2\end{pmatrix},\qquad C=C_1\otimes C_1.
 \label{eq:C}
\end{equation}
The eigenvalues of $C_1C_1^\dagger$ are
\begin{equation}
 \lambda_{\pm}=\frac{13\pm3\sqrt{17}}{26}.
 \label{eq:lambdapm}
\end{equation}
Consequently the four nonzero Schmidt weights at $W_0$ are
$\lambda_-^2,\lambda_-\lambda_+,\lambda_-\lambda_+,\lambda_+^2$, and
\begin{equation}
 F(W_0)=2h_2(\lambda_+)=0.329435221224950\ldots .
 \label{eq:F0}
\end{equation}
The one-copy optimality needed to identify this as twice the one-copy EoP is
the exact result of Ref.~\cite{NegariKhanian2026}; the local theorem below
does not reprove that global one-copy statement.

Local unitaries on the ancillary factors preserve the entropy.  Define their
image and orbit by
\begin{equation}
 \cG=\{U\otimes V:U,V\in U(4)\},\qquad \cO=\cG W_0.
 \label{eq:gauge}
\end{equation}

\begin{theorem}[minimal-split local stability]
\label{thm:main}
For the state $\rho_{1/2}$ in Eq.~\eqref{eq:state}, there is a neighborhood
$\mathcal N$ of $\cO$ in $U(16)$ such that
\begin{equation}
 F(W)\geq F(W_0)\quad (W\in\mathcal N),
\end{equation}
with equality only for $W\in\cO$ after $\mathcal N$ is reduced if necessary.
Thus $W_0$ is a strict local minimum modulo $\cG$ among all complex
orthonormal flags in the fixed $(4,4)$ ancillary split, including paths that
open new Schmidt support.
\end{theorem}

\section{Tangent geometry and gauge quotient}

Index a flag input by $x=(I,J)$ and a flag output by $y=(u,v)$, with all four
indices in $\{0,1,2,3\}$.  Across $AA':BB'$, the $16\times16$ coefficient
matrix of Eq.~\eqref{eq:flags} is
\begin{equation}
 M(W)_{(I,u),(J,v)}=C_{IJ}W_{(u,v),(I,J)},\qquad F(W)=S(MM^\dagger).
 \label{eq:Mmap}
\end{equation}
At $W_0$, only the rows $(I,I)$ and columns $(J,J)$ are active; their
$4\times4$ block is the invertible matrix $C$.

Let $K^\dagger=-K$ be a tangent and use $W(t)=\exp(tK)$.  If $P$ and $Q$
project onto the active row and column spaces, the normal--normal part of the
first variation is the real-linear map
\begin{equation}
 N(K)_{(I,u),(J,v)}=C_{IJ}K_{(u,v),(I,J)},\qquad
 u\neq I,\quad v\neq J.
 \label{eq:Nmap}
\end{equation}
Because every $C_{IJ}$ is nonzero, $N(K)=0$ precisely when no off-diagonal
entry of $K$ connects two points of the $4\times4$ flag grid whose two
coordinates both change.

This gives an exact count.  The kernel contains 16 imaginary diagonal
generators.  There are 24 unordered grid edges sharing $I$ and 24 sharing
$J$, each with real and imaginary anti-Hermitian generators.  Hence
\begin{equation}
 \dim_{\mathbb R}\ker N=16+2(24+24)=112,
 \qquad \dim_{\mathbb R}(\ker N)^\perp=144.
 \label{eq:split}
\end{equation}
The tangent to $\cO$ consists of $X\otimes I_4+I_4\otimes Y$ with
$X,Y\in\mathfrak u(4)$.  The kernel of $(X,Y)\mapsto X\otimes I+I\otimes Y$
is $(\ii cI,-\ii cI)$, so this gauge tangent has dimension $31$ and lies in
$\ker N$.

For completeness, an exact orthonormal complement of the gauge tangent in
$\ker N$ is constructed from the three Helmert contrasts
\begin{equation}
 h_1=\frac{(1,-1,0,0)}{\sqrt2},\quad
 h_2=\frac{(1,1,-2,0)}{\sqrt6},\quad
 h_3=\frac{(1,1,1,-3)}{\sqrt{12}}.
 \label{eq:helmert}
\end{equation}
It contains nine diagonal phase tables $h_p h_q^{\mathsf T}$, 36
same-$I$ rook-edge contrasts, and 36 same-$J$ rook-edge contrasts.  Thus the
physical fixed-rank quotient has dimension $9+36+36=81$.  Figure~\ref{fig:split}
summarizes the decomposition.

\input{manuscript/figure-tangent.tex}

The 144 normalized opening generators have mutually orthogonal images under
$N$.  For an unordered edge between $x=(I,J)$ and $y=(u,v)$ with both
coordinates different, either its real or imaginary normalized generator has
\begin{equation}
 \lVert N(K)\rVert_F^2=\frac{C_{IJ}^2+C_{uv}^2}{2}.
 \label{eq:opening-scale}
\end{equation}
Exact enumeration gives
\begin{equation}
 \min_{K\in(\ker N)^\perp,\,\lVert K\rVert_F=1}
 \lVert N(K)\rVert_F^2=\frac{5}{338}.
 \label{eq:opening-gap}
\end{equation}
In particular, the opening map is quantitatively injective on its
144-dimensional complement.

The point $W_0$ is stationary in every tangent direction.  For a diagonal
$K_{(I,J),(I,J)}=\ii\theta_{IJ}$, the active part of
$\rho'(0)$, where $\rho=MM^\dagger$, is purely imaginary and antisymmetric
after removal of the factor $\ii$, whereas $\log\rho(0)$ is real symmetric.
Their trace pairing vanishes.  Rook-edge directions have only
active--normal first variations, and the remaining directions have
normal--normal first variations, so they also have zero first entropy
variation.  In an opening direction the new weights are $O(t^2)$ and their
entropy is $O(-t^2\log t^2)$, whose first derivative at zero is again zero.

\section{Certified Hessian on the fixed-rank quotient}

We equip $\mathfrak u(16)$ with
$\langle K,L\rangle=\operatorname{Re}\Tr(K^\dagger L)$.  Write
\begin{equation}
 M(t)=M_0+tM_1+\frac{t^2}{2}M_2+O(t^3),
 \quad M_1=L(K),\quad M_2=L(K^2),
\end{equation}
where $L$ is the linear coefficient map in Eq.~\eqref{eq:Mmap}.  Then
\begin{align}
 \rho_1&=M_1M_0^\dagger+M_0M_1^\dagger,\\
 \rho_2&=M_2M_0^\dagger+M_0M_2^\dagger+2M_1M_1^\dagger.
 \label{eq:rho12}
\end{align}
Let $P$ now denote the four-dimensional positive support of $\rho_0$ and
$Q_0=I-P$.  Let $U$ diagonalize $P\rho_0P$, with eigenvalues
$\{\lambda_a\}_{a=1}^4=\{\lambda_-^2,\lambda_-\lambda_+,
\lambda_-\lambda_+,\lambda_+^2\}$.  Put
\begin{equation}
 X=U^\dagger P\rho_1PU,\quad Z=U^\dagger P\rho_2PU,
 \quad Y=Q_0\rho_1PU.
\end{equation}
For $f(x)=-x\log_2x$, define the divided difference
\begin{equation}
 f'^{[1]}(x,y)=
 \begin{cases}
 [f'(x)-f'(y)]/(x-y),&x\neq y,\\
 f''(x),&x=y.
 \end{cases}
\end{equation}
Hermitian spectral perturbation, including the double eigenvalue
$\lambda_-\lambda_+$, gives the second directional derivative
\begin{equation}
 q(K)=\sum_a f'(\lambda_a)Z_{aa}
 +\sum_{a,b}f'^{[1]}(\lambda_a,\lambda_b)|X_{ab}|^2
 +2\sum_a\frac{f'(\lambda_a)}{\lambda_a}\lVert Y_{:a}\rVert_2^2.
 \label{eq:hessian}
\end{equation}
The last term records motion of the positive eigenspace into the zero
eigenspace.  If $K\in\ker N$, the small singular values of $M(\exp(tK))$ are
$O(t^2)$, so their squared Schmidt weights are $O(t^4)$ and their entropy is
$O(-t^4\log|t|)$.  They therefore have zero second derivative.  Thus
Eq.~\eqref{eq:hessian} is the Hessian of the full entropy on $\ker N$.

The 81-by-81 matrix $H$ of $q$ in the basis described after
Eq.~\eqref{eq:helmert} was evaluated with 320-bit Arb midpoint-radius balls
through \texttt{python-flint}; Arb implements outward-rounded ball
arithmetic~\cite{Johansson2017}.  The calculation uses exact rational and
radical input expressions and analytic divided differences, not finite
differences.  It then performs an interval, no-pivot $LDL^{\mathsf T}$
factorization of
\begin{equation}
 H-\frac{159}{1000}I_{81}.
 \label{eq:shift}
\end{equation}
Every interval pivot is strictly positive.  Sylvester's criterion therefore
certifies
\begin{equation}
 \boxed{\lambda_{\min}(H)>0.159}
 \quad\text{bits per squared unit tangent.}
 \label{eq:hessian-gap}
\end{equation}
Table~\ref{tab:certificate} records the decisive checks.  The smallest pivot
is not itself an eigenvalue bound; it certifies the factorization of the
already shifted matrix.  Likewise, the binary64 midpoint eigenvalue is a
diagnostic only.

\begin{table}[t]
\centering
\caption{Interval-certificate summary.  All inequalities in the certificate
use outward-rounded balls; the midpoint eigensolver is not part of the proof.}
\label{tab:certificate}
\begin{tabular}{@{}ll@{}}
\toprule
Quantity & Certified or recorded value \\
\midrule
Quotient dimension & $81=9+36+36$ \\
Arithmetic precision & 320 bits \\
Largest Hessian-entry radius & $<1.64\times10^{-93}$ \\
Rational shift & $159/1000$ \\
Strictly positive $LDL^{\mathsf T}$ pivots & 81 of 81 \\
Smallest pivot lower endpoint & $>0.0290093$ \\
Reconstruction residual & every entry encloses zero \\
Opening-map least squared scale & $5/338$ exactly \\
Midpoint $\lambda_{\min}$ (diagnostic) & $0.159148973002032$ \\
Independent precision replays & 256 and 384 bits, both pass \\
\bottomrule
\end{tabular}
\end{table}

\section{Uniform control of support opening}

Positive definiteness on $\ker N$ does not by itself prove
Theorem~\ref{thm:main}: the reduced state at $W_0$ has twelve zero
eigenvalues.  We now keep their logarithmic entropy explicitly.

Use the active and normal row and column decompositions to write, near $W_0$,
\begin{equation}
 M(W)=\begin{pmatrix}A&B\\E&D\end{pmatrix},\qquad A(W_0)=C,
\end{equation}
and define the analytic Schur complement
\begin{equation}
 \mathcal Z(W)=D-EA^{-1}B.
 \label{eq:schur}
\end{equation}
Its derivative at $W_0$ is $N$.  The exact factorization
\begin{equation}
 M=
 \begin{pmatrix}I&0\\EA^{-1}&I\end{pmatrix}
 \begin{pmatrix}A&0\\0&\mathcal Z\end{pmatrix}
 \begin{pmatrix}I&A^{-1}B\\0&I\end{pmatrix}
 \label{eq:blockfact}
\end{equation}
has uniformly bounded outer factors and inverses in a sufficiently small
neighborhood.  By the characterization of the Frobenius distance to matrices
of rank at most four, the total small-weight mass
\begin{equation}
 \tau(W)=\sum_{j=5}^{16}s_j(M(W))^2
 \label{eq:tau}
\end{equation}
satisfies $\tau(W)\geq c_0\lVert\mathcal Z(W)\rVert_F^2$ there for some
$c_0>0$.

Let $\Pi$ be the orthogonal projection onto the real 144-dimensional image of
$N$ and set $z=\Pi\mathcal Z(W)$.  Work on a local slice perpendicular to the
31 gauge directions.  The exact gap in Eq.~\eqref{eq:opening-gap} makes the
derivative from opening coordinates to $z$ invertible.  The inverse-function
theorem therefore gives local coordinates $(k,z)$ on the slice, where
$k\in\mathbb R^{81}$ is the fixed-rank quotient coordinate, and after
shrinking the chart,
\begin{equation}
 \tau(k,z)\geq c\lVert z\rVert^2.
 \label{eq:taulower}
\end{equation}

Let $G(k,z)$ be the entropy contribution of the four eigenvalues separated
from zero.  This is analytic despite their internal double eigenvalue.  Full
stationarity and Eq.~\eqref{eq:hessian-gap}, followed by an ordinary Taylor
bound on the remaining Hessian blocks, yield constants $A,B>0$ such that
\begin{equation}
 G(k,z)-G(0,0)\geq A\lVert k\rVert^2
 -B\lVert k\rVert\lVert z\rVert-B\lVert z\rVert^2.
 \label{eq:Gbound}
\end{equation}
For sufficiently small $\tau$, the entropy $T$ of the other twelve weights
obeys
\begin{equation}
 T=\sum_{j=5}^{16}-s_j^2\log_2s_j^2
 \geq-\tau\log_2\tau
 \geq-c\lVert z\rVert^2\log_2(c\lVert z\rVert^2).
 \label{eq:tailbound}
\end{equation}
The first inequality uses $s_j^2\leq\tau$; the second uses
Eq.~\eqref{eq:taulower} and monotonicity of $-x\log_2x$ near zero.

Young's inequality absorbs the mixed term in Eq.~\eqref{eq:Gbound} into half
of the positive $k$ term plus a constant multiple of $\lVert z\rVert^2$.
The logarithm in Eq.~\eqref{eq:tailbound} uniformly dominates every such
bounded quadratic penalty as $z\to0$.  Hence $F(k,z)>F(0,0)$ at every
sufficiently small nonzero point of the slice.  Equality occurs only at its
origin.  Finally, the compact ancillary action admits a tubular local slice,
and $F$ is constant on each ancillary orbit.  This transfers the strict slice
minimum to a neighborhood of $\cO$ and completes the proof of
Theorem~\ref{thm:main}.

\section{Interpretation and limitations}

The result controls all complex flags in the smallest bipartite ancillary
space capable of holding the sixteen occupied two-copy branches.  It is
stronger than a scan of sign patterns or a direction-by-direction numerical
test: the fixed-rank part is certified on the complete gauge quotient, and
the rank boundary is controlled uniformly rather than inferred from an
ordinary Hessian.

At the same time, the statement is intentionally narrow.  The neighborhood
radius is existential and has not been optimized.  Only $r=1/2$ and
$\alpha=1$ are treated.  A purification at finite distance in $U(16)$ could
still have smaller entropy, and a larger ancillary split has additional
directions absent here.  Thus neither
$E_P(\rho_{1/2}^{\otimes2})=2E_P(\rho_{1/2})$ nor any regularized correlation
cost follows.  The rigorous consequence is instead a landscape constraint:
at this state and in this split, any von Neumann improvement must occupy a
separate finite-distance basin or enter through ancillary enlargement.

Two natural next steps are to turn the existence proof into an explicit
certified neighborhood radius, and to repeat the tangent-plus-boundary
analysis in a padded ancillary embedding.  Either extension requires new
scientific certification and is not assumed here.

\section{Reproducibility and data availability}

The accompanying source package contains the exact basis construction, the
analytic Hessian evaluator, the 320-bit Arb certificate, its 256- and 384-bit
precision replays, independent finite-difference diagnostics, and the source
baseline reconstruction.  The designated certificate is generated from the
project root by
\begin{verbatim}
python3 evidence/certify_u16_hessian.py
\end{verbatim}
and must report 81 positive pivots for the shift in
Eq.~\eqref{eq:shift}.  The recorded environment uses Python 3.12.14,
\texttt{python-flint} 0.9.0, NumPy 2.3.5, SciPy 1.18.1, and mpmath 1.3.0.
The reproducibility README distinguishes the rigorous interval output from
diagnostic floating-point calculations.  No experimental or personal data
were used.  A permanent public repository and identifier remain decisions
for the author before submission.

\section*{Declarations}

\paragraph{AI use.}
Generative AI (OpenAI Codex, GPT-5.6 Sol, accessed September 2026) assisted
with formulating the scoped question, mathematical derivations, code,
literature checks, computational verification, an independent AI referee
replay, and manuscript drafting.  AI is not an author and cannot assume
accountability.  The named author must independently review the claims,
proof, code, citations, and prose, and approve an accurate final declaration
before any submission.

\paragraph{Author contributions, funding, and competing interests.}
These declarations require confirmation and completion by the author before
submission.  No external submission has been made from this draft.

\bibliographystyle{unsrtnat}
\bibliography{manuscript/references}

\end{document}
